\chapter{Introduction}
\label{chap:main-introduction}
\label{chap:introduction}

Many data sets record curves rather than isolated measurements: growth
trajectories, electricity demand, biomedical signals, and spectra are familiar
examples. Their measurement locations have an order and a meaningful
relationship to neighbouring locations. Functional data analysis uses this
structure to describe variation in whole trajectories and relate it to other
variables \parencite{ramsay_silverman_2005}. This thesis studies the
scalar-on-function linear model, in which a response depends on the integral
of a predictor curve against an unknown slope function.

Estimating that slope is an inverse problem. The predictor covariance
operator determines which slope directions the data identify and how much
information they provide about each direction. In an infinite-rank model,
its positive eigenvalues approach zero. Inverting small eigenvalues can
greatly amplify sampling errors, even when the slope is uniquely identified
on the parameter space of interest. A finite representation makes the
calculation possible but does not remove the underlying instability
\parencite{cardot_ferraty_sarda_1999,hall_horowitz_2007}.

Regularisation controls this amplification by restricting what the estimator
attempts to recover. Functional principal component regression (FPCR)
retains a leading empirical eigenspace. Functional partial least squares
(FPLS) constructs directions using the relationship between predictor and
response, rather than predictor variation alone
\parencite{preda_saporta_2005,delaigle_hall_2012}. Its Krylov formulation
connects functional regression to iterative methods for linear systems.
The number of admitted directions or iterations becomes a regularisation
parameter: it determines how much of the empirical inverse problem is solved.

There are two reasons why the same number of components need not work for
every purpose. First, slope estimation and prediction assign different
importance to errors. A slope error in a direction with little predictor
variation can contribute substantially to integrated squared error while
having little effect on predictions \parencite{cai_hall_2006}. Inference
adds a third criterion. For the moment-based test considered here, a
sufficient approximation condition requires the remaining fitting residual
to be negligible relative to sampling uncertainty. Stopping that is useful
for regularised estimation may leave too much residual for this argument.

Second, a well-defined mathematical search space can be difficult to
represent numerically. Successive powers of an empirical covariance operator
may become nearly parallel. Solving least squares in these raw coordinates
can then be unstable. Arnoldi orthogonalisation represents the same Krylov
space through an orthonormal basis, while iterative recurrences avoid forming
the complete power-basis equations. Statistical regularisation and numerical
reliability must therefore be examined together.

\section{Questions and contribution}

The thesis compares four complete procedures: FPCR, raw power-basis FPLS,
Arnoldi FPLS with reorthogonalisation, and CG--FPLS with discrepancy stopping.
The last name follows the underlying study of
\textcite{babii_carrasco_tsafack_2026}; its implemented updates minimise
the moment residual and are conjugate-residual updates in standard numerical
terminology. They generally differ from the response least-squares target
used by Raw and Arnoldi before convergence. This distinction is maintained
throughout the comparison.

Four questions guide the analysis:
\begin{enumerate}
    \item How do spectral and iterative regularisation compare for slope
    recovery and independent prediction when the slope and covariance change?
    \item How much does the numerical representation affect theoretically
    equivalent raw and Arnoldi fits?
    \item Why can a stopping index adequate for estimation be inadequate
    for the stated inference procedure?
    \item Which of these effects remain visible in a real spectroscopic
    prediction problem?
\end{enumerate}

The contribution is methodological, computational, and expository. It brings
these questions into a common analysis, rather than proposing a new
asymptotic theory. Existing functional-regression results motivate the
comparison, and the supporting proofs make its operator identities and
inferential assumptions explicit. Comparative work already shows that the
relative merits of FPCR and FPLS depend on covariance, slope alignment, and
tuning \parencite{febrero_bande_et_al_2017}. Here the additional focus is on
how numerical implementation and the statistical target affect that comparison.

The estimation experiment uses $5{,}000$ simulated training samples in each
of three controlled models. All four methods receive the same data, but
retain their own fitting and tuning rules. A separate CG--FPLS study
examines the remaining moment residual, null rejection probabilities, power,
and confidence sets. Finally, an application predicts pharmaceutical-tablet
assay from near-infrared spectra, using $195$ development observations and
a designated $460$-tablet benchmark. The application concerns prediction
and numerical stability; its unknown physical slope does not permit a
direct assessment of slope recovery.

\section{Organisation of the thesis}

Chapter~\ref{chap:main-methods} gives the model, essential operator
background, and four procedures. Chapter~\ref{chap:main-simulation}
combines the simulation design with its principal estimation and prediction
results. Chapter~\ref{chap:main-inference} develops the moment-based
testing argument and evaluates its finite-sample behaviour.
Chapter~\ref{chap:main-empirical} presents the spectroscopy application,
and Chapter~\ref{chap:main-conclusion} draws together the implications
and limitations.

The appendices retain the extended work behind this argument. They contain
the functional-analysis background, mathematical proofs, complete algorithmic
and experimental specifications, detailed results, and additional diagnostic
figures. The main text presents the definitions and evidence needed to
follow the conclusions; the appendices allow their assumptions and numerical
details to be examined more closely.
